\documentclass[11pt,a4paper]{article}
% preamble.tex — estilo común de todos los apuntes Froth.
% Cada apunte hace % preamble.tex — estilo común de todos los apuntes Froth.
% Cada apunte hace % preamble.tex — estilo común de todos los apuntes Froth.
% Cada apunte hace \input{preamble} justo después de \documentclass.
\usepackage[margin=2.4cm]{geometry}
\usepackage{xcolor}
\definecolor{litblue}{RGB}{31,79,121}
\definecolor{litgreen}{RGB}{27,94,32}
\definecolor{litorange}{RGB}{230,81,0}
\usepackage{parskip}
\usepackage{titlesec}
\titleformat{\section}{\Large\bfseries\color{litblue}}{\thesection}{0.6em}{}
\titleformat{\subsection}{\large\bfseries\color{litblue}}{\thesubsection}{0.6em}{}
\usepackage{tcolorbox}
\usepackage{fancyhdr}
\pagestyle{fancy}
\fancyhf{}
\fancyhead[L]{\small\color{litblue}Froth — Apuntes de ML}
\fancyhead[R]{\small\thepage}
\renewcommand{\headrulewidth}{0.4pt}
\usepackage[colorlinks=true,linkcolor=litblue,urlcolor=litblue]{hyperref}

% Cabecera de cada apunte: \cabecera{numero}{titulo}
\newcommand{\cabecera}[2]{%
  \begin{tcolorbox}[colback=litblue,coltext=white,colframe=litblue,arc=2mm]
    {\Large\bfseries Apunte #1 — #2}\\[3pt]
    {\small Proyecto Froth · aprender Machine Learning construyendo}
  \end{tcolorbox}\vspace{0.8em}}

% Cajas temáticas
\newtcolorbox{concepto}[1]{colback=blue!4,colframe=litblue,title={Concepto: #1},arc=1.5mm}
\newtcolorbox{practica}{colback=green!5,colframe=litgreen,title={Pruébalo tú (teclado en mano)},arc=1.5mm}
\newtcolorbox{ojo}{colback=orange!8,colframe=litorange,title={Ojo / error típico},arc=1.5mm}
\newtcolorbox{resumen}{colback=gray!8,colframe=gray!60,title={En una frase},arc=1.5mm}

\newcommand{\codigo}[1]{\texttt{#1}}
 justo después de \documentclass.
\usepackage[margin=2.4cm]{geometry}
\usepackage{xcolor}
\definecolor{litblue}{RGB}{31,79,121}
\definecolor{litgreen}{RGB}{27,94,32}
\definecolor{litorange}{RGB}{230,81,0}
\usepackage{parskip}
\usepackage{titlesec}
\titleformat{\section}{\Large\bfseries\color{litblue}}{\thesection}{0.6em}{}
\titleformat{\subsection}{\large\bfseries\color{litblue}}{\thesubsection}{0.6em}{}
\usepackage{tcolorbox}
\usepackage{fancyhdr}
\pagestyle{fancy}
\fancyhf{}
\fancyhead[L]{\small\color{litblue}Froth — Apuntes de ML}
\fancyhead[R]{\small\thepage}
\renewcommand{\headrulewidth}{0.4pt}
\usepackage[colorlinks=true,linkcolor=litblue,urlcolor=litblue]{hyperref}

% Cabecera de cada apunte: \cabecera{numero}{titulo}
\newcommand{\cabecera}[2]{%
  \begin{tcolorbox}[colback=litblue,coltext=white,colframe=litblue,arc=2mm]
    {\Large\bfseries Apunte #1 — #2}\\[3pt]
    {\small Proyecto Froth · aprender Machine Learning construyendo}
  \end{tcolorbox}\vspace{0.8em}}

% Cajas temáticas
\newtcolorbox{concepto}[1]{colback=blue!4,colframe=litblue,title={Concepto: #1},arc=1.5mm}
\newtcolorbox{practica}{colback=green!5,colframe=litgreen,title={Pruébalo tú (teclado en mano)},arc=1.5mm}
\newtcolorbox{ojo}{colback=orange!8,colframe=litorange,title={Ojo / error típico},arc=1.5mm}
\newtcolorbox{resumen}{colback=gray!8,colframe=gray!60,title={En una frase},arc=1.5mm}

\newcommand{\codigo}[1]{\texttt{#1}}
 justo después de \documentclass.
\usepackage[margin=2.4cm]{geometry}
\usepackage{xcolor}
\definecolor{litblue}{RGB}{31,79,121}
\definecolor{litgreen}{RGB}{27,94,32}
\definecolor{litorange}{RGB}{230,81,0}
\usepackage{parskip}
\usepackage{titlesec}
\titleformat{\section}{\Large\bfseries\color{litblue}}{\thesection}{0.6em}{}
\titleformat{\subsection}{\large\bfseries\color{litblue}}{\thesubsection}{0.6em}{}
\usepackage{tcolorbox}
\usepackage{fancyhdr}
\pagestyle{fancy}
\fancyhf{}
\fancyhead[L]{\small\color{litblue}Froth — Apuntes de ML}
\fancyhead[R]{\small\thepage}
\renewcommand{\headrulewidth}{0.4pt}
\usepackage[colorlinks=true,linkcolor=litblue,urlcolor=litblue]{hyperref}

% Cabecera de cada apunte: \cabecera{numero}{titulo}
\newcommand{\cabecera}[2]{%
  \begin{tcolorbox}[colback=litblue,coltext=white,colframe=litblue,arc=2mm]
    {\Large\bfseries Apunte #1 — #2}\\[3pt]
    {\small Proyecto Froth · aprender Machine Learning construyendo}
  \end{tcolorbox}\vspace{0.8em}}

% Cajas temáticas
\newtcolorbox{concepto}[1]{colback=blue!4,colframe=litblue,title={Concepto: #1},arc=1.5mm}
\newtcolorbox{practica}{colback=green!5,colframe=litgreen,title={Pruébalo tú (teclado en mano)},arc=1.5mm}
\newtcolorbox{ojo}{colback=orange!8,colframe=litorange,title={Ojo / error típico},arc=1.5mm}
\newtcolorbox{resumen}{colback=gray!8,colframe=gray!60,title={En una frase},arc=1.5mm}

\newcommand{\codigo}[1]{\texttt{#1}}

\begin{document}
\cabecera{22}{Las vistas finales: red con física y circle packing}

\section{Tres vistas, tres preguntas}

Tu app ahora tiene \textbf{cuatro pestañas}, y cada vista responde una pregunta distinta:

\begin{center}
\begin{tabular}{lp{5.2cm}p{4.6cm}}
\textbf{Vista} & \textbf{Responde a...} & \textbf{Posición significa...} \\
\hline
Topic map (scatter) & ¿Qué grupos hay y qué tan cerca están? & Distancia semántica (UMAP) \\
Network (red) & ¿Quién se conecta con quién? ¿Hubs? ¿Puentes? & La física acomoda; los enlaces son lo real \\
Packed bubbles & ¿Cómo se reparte el campo? (presentar) & \textbf{Nada} — solo pertenencia y tamaño \\
\end{tabular}
\end{center}

Saber \emph{qué significa la posición en cada vista} te salva de sobre-interpretar: en el
scatter la cercanía ES información; en el packing es solo estética.

\section{La red con física (pyvis)}

\begin{concepto}{layout force-directed (dirigido por fuerzas)}
La pestaña Network simula \textbf{física}: cada enlace es un \emph{resorte} que atrae a los
papers conectados, y todos los nodos se \emph{repelen} un poco entre sí. El sistema se suelta
y busca equilibrio solo: las comunidades muy conectadas quedan juntas y compactas, los
puentes quedan estirados en medio. Por eso al \textbf{arrastrar un nodo} toda la red
reacciona — esa es la ``animación'': no es decorado, es el equilibrio recalculándose.
La librería: \codigo{pyvis} (envuelve vis.js, un motor JavaScript).
\end{concepto}

\section{El circle packing (circlify)}

\begin{concepto}{empaquetado de círculos}
La vista Packed bubbles coloca un \textbf{círculo grande por subtema} y dentro empaqueta un
círculo por paper (área $=$ citas), \textbf{sin solapes}. Calcular dónde va cada círculo para
que quepan apretados es un problema geométrico clásico; la librería \codigo{circlify} lo
resuelve y nosotros solo dibujamos (con Plotly, para conservar el hover). Es la vista
``de presentación'' — la del gráfico de packaging de la UE que inspiró el proyecto.
\end{concepto}

\begin{ojo}
En el packing, dos papers vecinos NO son más parecidos que dos lejanos: la posición dentro
del círculo la decide la geometría del empaquetado, no el significado. Membresía (qué círculo
grande) y tamaño (citas) son la única información. Elegante para comunicar; no para analizar.
\end{ojo}

\section{Dónde vive todo}

\begin{itemize}
  \item \codigo{froth/visualize.py}: \codigo{network\_html()} (red pyvis autocontenida),
        \codigo{packed\_figure()} (packing con circlify+Plotly), y sus versiones
        \codigo{save\_*}/\codigo{plot\_*} que escriben HTML en \codigo{3\_Resultados/}.
  \item \codigo{app.py}: pestañas \textbf{Network} y \textbf{Packed bubbles} nuevas. La red
        se cachea por versión de datos (mismo truco del Apunte del slider).
  \item Archivos compartibles: \codigo{topic\_map.html}, \codigo{paper\_network.html},
        \codigo{packed\_bubbles.html} — doble clic y se abren en cualquier navegador.
\end{itemize}

\begin{resumen}
Tres vistas para tres preguntas: scatter (distancias semánticas), red con física (conexiones,
hubs, puentes — arrastrable porque es un sistema de resortes) y circle packing (presentación
pura, posición sin significado). Todas en la app y como HTMLs compartibles.
\end{resumen}

\end{document}
